\documentclass[11pt]{article}

\usepackage[margin=1in]{geometry}
\usepackage{amsmath}
\usepackage{amssymb}
\usepackage{graphicx}
\usepackage{booktabs}
\usepackage{hyperref}
\usepackage{xcolor}
\usepackage{listings}
\usepackage{enumitem}

\hypersetup{
  colorlinks=true,
  linkcolor=blue,
  urlcolor=blue,
  citecolor=blue
}

\lstdefinestyle{code}{
  basicstyle=\ttfamily\small,
  breaklines=true,
  frame=single,
  columns=fullflexible,
  keepspaces=true,
  showstringspaces=false
}
\lstset{style=code}

\title{CyberRunner Reinforcement Learning Training Report}
\author{Trung Bao}
\date{\today}

\begin{document}
\maketitle

\begin{abstract}
This report explains the CyberRunner training system from the ground up. It is written for readers who may not know robotics, reinforcement learning, or DreamerV3 before reading it. The report first introduces the basic ideas of reinforcement learning, then explains how those ideas are used on the physical CyberRunner maze robot. It covers the robot hardware, observations, actions, reward, replay data, checkpoints, policy learning, DreamerV3, and the practical workflow used to keep the best policies instead of only the latest one.
\end{abstract}

\tableofcontents
\newpage

\section{Introduction}

CyberRunner is a robotic maze system. A small ball moves on top of a board, and the board can tilt in two directions. By tilting the board correctly, the robot can guide the ball along a path and toward the end of the maze.

The main goal is:

\begin{center}
\fbox{\parbox{0.88\linewidth}{
\centering
learn a control policy that moves the ball through the maze reliably, using camera feedback and servo commands
}}
\end{center}

This is difficult because the robot is physical. The ball moves quickly, the camera can be noisy, the servo response has delay, and a small wrong motion can make the ball fall off the path.

\section{Basic Reinforcement Learning Concepts}

Reinforcement learning, or RL, is a way for a computer program to learn by trying actions and receiving rewards. It is different from supervised learning. In supervised learning, the program is given correct answers. In reinforcement learning, the program is not told the correct action directly. Instead, it must discover useful behavior by trial and error.

\subsection{Agent}

The \textbf{agent} is the learner. In this project, the agent is the DreamerV3 neural network system running on the training server. It looks at the current robot observation and chooses what motor command to send.

\subsection{Environment}

The \textbf{environment} is the outside world that the agent interacts with. In CyberRunner, the environment includes:

\begin{itemize}[leftmargin=*]
  \item the physical maze board,
  \item the ball,
  \item the camera and state estimation,
  \item the servos,
  \item the ROS and TCP software that sends data between the robot and server.
\end{itemize}

\subsection{Observation}

An \textbf{observation} is what the agent can see at one moment. It is not the whole real world, only the information given to the policy. For CyberRunner, this includes a small camera image, ball position, board angle information, and path progress information.

\subsection{Action}

An \textbf{action} is what the agent decides to do. In CyberRunner, the action is converted into two servo velocity commands, one for each board axis.

\subsection{Reward}

A \textbf{reward} is a number that tells the agent whether the last behavior was useful. A positive reward means the robot made progress. A low or zero reward means it did not make useful progress.

\subsection{Policy}

The \textbf{policy} is the part of the agent that chooses actions. In simple terms:

\begin{lstlisting}
observation -> policy -> action
\end{lstlisting}

Mathematically:

\begin{equation}
\pi(o_t) = a_t
\end{equation}

where $o_t$ is the observation at time $t$, $a_t$ is the action, and $\pi$ is the policy.

\subsection{Episode}

An \textbf{episode} is one attempt at the maze. It starts when the robot resets and the ball is detected. It ends when the ball is lost, the ball goes off the path, the goal is reached, or the maximum number of steps is reached.

\subsection{Checkpoint}

A \textbf{checkpoint} is a saved copy of the trained neural network. It lets training continue later or lets the robot run a saved policy. A checkpoint is important because the latest policy is not always the best policy.

\section{CyberRunner System Overview}

The project uses a split system:

\begin{itemize}[leftmargin=*]
  \item The server runs DreamerV3 training.
  \item The local PC runs the camera, state estimation, servo control, and TCP bridge.
  \item The server sends actions to the local PC.
  \item The local PC sends observations and rewards back to the server.
\end{itemize}

The main TCP training environment is:

\begin{lstlisting}
cyberrunner_dreamer/cyberrunner_dreamer/env_tcp.py
\end{lstlisting}

The local ROS environment is:

\begin{lstlisting}
cyberrunner_dreamer/cyberrunner_dreamer/env.py
\end{lstlisting}

The control loop is:

\begin{center}
\fbox{\parbox{0.92\linewidth}{
\centering
camera image and state estimation $\rightarrow$ observation $\rightarrow$ Dreamer policy
$\rightarrow$ velocity action $\rightarrow$ servo motion $\rightarrow$ new ball position
}}
\end{center}

\section{What the Robot Learns}

The robot does not learn the maze geometry by manually programmed rules. Instead, the policy learns how to choose motor commands based on observations.

The trained behavior should answer questions like:

\begin{itemize}[leftmargin=*]
  \item If the ball is left of the path, how should the board tilt?
  \item If the ball is moving too fast, how should the robot slow it down?
  \item If the ball is near a turn, how should the robot prepare?
  \item If the ball starts drifting toward failure, how can the robot recover?
\end{itemize}

This is why the policy is the main trained object.

\section{Our Policy in This Project}

In this project, the word \textbf{policy} refers to the neural network inside DreamerV3 that outputs the robot action. It is not a hand-written controller. It is a learned actor network.

The policy is used every control step:

\begin{lstlisting}
camera/state data -> Dreamer encoder -> latent state -> actor policy -> action -> servo velocity
\end{lstlisting}

The policy does not directly look only at the raw camera image. First, DreamerV3 converts the observation into a hidden representation called a \textbf{latent state}. The policy then uses this latent state to choose an action.

\subsection{Policy Input}

The policy input is the Dreamer latent state. In the current DreamerV3 configuration, the actor uses:

\begin{lstlisting}
actor.inputs: [deter, stoch]
\end{lstlisting}

These two parts come from the recurrent state-space model:

\begin{itemize}[leftmargin=*]
  \item \texttt{deter}: a deterministic memory state. This helps remember recent history.
  \item \texttt{stoch}: a stochastic latent state. This represents uncertainty and hidden state information.
\end{itemize}

This matters because the robot cannot be controlled well from only one image. The ball velocity and recent motion are important. The recurrent latent state gives the policy memory of what just happened.

\subsection{Policy Network Architecture}

The actor policy is an MLP, which means a multi-layer perceptron. In the current DreamerV3 configuration, it uses:

\begin{lstlisting}
actor: layers=5, units=1024, act=silu, norm=layer
actor_dist_cont: normal
actor_opt: adam, lr=3e-5
\end{lstlisting}

In plain language:

\begin{itemize}[leftmargin=*]
  \item The policy is a neural network with 5 hidden layers.
  \item Each hidden layer has 1024 units.
  \item It uses the SiLU activation function.
  \item It uses layer normalization to make training more stable.
  \item It is trained with the Adam optimizer.
\end{itemize}

\subsection{Policy Output}

CyberRunner has two continuous action dimensions:

\begin{equation}
a_t = [a_{1,t}, a_{2,t}]
\end{equation}

The action space is:

\begin{lstlisting}[language=Python]
gym.spaces.Box(-1.0, 1.0, (2,))
\end{lstlisting}

So each action value should be between $-1$ and $1$. These two numbers are later scaled into two servo velocity commands.

\subsection{The Policy Outputs a Distribution}

The actor does not output only one fixed action number. It outputs a \textbf{normal distribution} for each action dimension. A normal distribution has:

\begin{itemize}[leftmargin=*]
  \item a mean, which is the center action the policy wants,
  \item a standard deviation, which controls how much randomness is used.
\end{itemize}

The DreamerV3 code uses:

\begin{lstlisting}
actor_dist_cont: normal
actor.minstd: 0.1
actor.maxstd: 1.0
\end{lstlisting}

This means the policy can still explore instead of always choosing the same exact action. Exploration is useful during training because the robot needs to try different behaviors to discover better ones.

The actor output mean is passed through a \texttt{tanh} function. This helps keep the action centered inside the valid range from $-1$ to $1$.

\subsection{Training Mode and Evaluation Mode}

During training, the policy samples actions from its action distribution. Sampling means the action has some randomness. This helps exploration.

During evaluation, the code uses the task policy instead of the exploration policy. In this repository, the evaluation path still samples from the task policy distribution, but it does not use the separate exploration behavior. This means evaluation is less exploratory than training, but it can still have some randomness because the policy distribution has standard deviation.

This is one reason one checkpoint can sometimes finish the map and sometimes fail. The physical system is noisy, and the policy action can still vary slightly.

\subsection{How the Policy Learns}

The policy is trained with an actor-critic method inside DreamerV3.

The two main learned parts are:

\begin{itemize}[leftmargin=*]
  \item \textbf{Actor}: chooses actions.
  \item \textbf{Critic}: estimates how good a state is.
\end{itemize}

The critic predicts future return. The actor is updated so that its actions lead to higher predicted return.

DreamerV3 does this using imagined rollouts. The world model imagines what may happen for several future steps. In the current configuration:

\begin{lstlisting}
imag_horizon: 15
\end{lstlisting}

That means the actor is trained using short imagined futures of about 15 latent steps. The actor learns from those imagined futures instead of waiting for every update to come directly from the physical robot.

\subsection{How the Policy Becomes Motor Motion}

The actor outputs two normalized values:

\begin{equation}
[a_1, a_2], \quad a_1,a_2 \in [-1,1]
\end{equation}

The environment converts them into servo velocity:

\begin{lstlisting}[language=Python]
action *= self.max_angle_vel
vel_1 = self.alpha_fac * action[0]
vel_2 = self.beta_fac * action[1]
\end{lstlisting}

For the restored Hiwonder/original setup:

\begin{lstlisting}[language=Python]
self.max_angle_vel = 80
self.alpha_fac = -1.0
self.beta_fac = -1.0
\end{lstlisting}

So if the policy outputs:

\begin{equation}
[0.5, -1.0]
\end{equation}

then the approximate command becomes:

\begin{equation}
[-40, 80]
\end{equation}

because each action is multiplied by 80 and then flipped by the negative sign factors.

\subsection{Why the Policy Can Fail Even After Learning}

Even a good policy can fail sometimes. This does not always mean the policy is useless. It can happen because:

\begin{itemize}[leftmargin=*]
  \item the camera measurement changes slightly,
  \item the ball starts with a different velocity,
  \item the servo response is delayed,
  \item the policy samples a slightly different action,
  \item the ball reaches a difficult part of the maze,
  \item the learned world model is not perfect.
\end{itemize}

For this reason, a good checkpoint should be evaluated over multiple attempts. The best policy is the one that succeeds reliably, not only the one that gets one lucky high score.

\section{Observation Details}

The observation contains several pieces of information:

\begin{itemize}[leftmargin=*]
  \item A small grayscale image of the board and ball, usually $64 \times 64 \times 1$.
  \item Board angles such as \texttt{alpha} and \texttt{beta}.
  \item Ball position such as \texttt{x\_b} and \texttt{y\_b}.
  \item Path information from the saved maze path.
  \item Progress and reward values used for logging and training.
\end{itemize}

The image gives visual information. The state values give compact physical information. The path progress tells the environment whether the ball is moving forward through the maze.

\section{Action and Servo Commands}

DreamerV3 outputs a continuous action:

\begin{equation}
a_t = [a_{1,t}, a_{2,t}], \quad a_{i,t} \in [-1, 1]
\end{equation}

This action is scaled into servo velocity commands. In the restored Hiwonder/original setup, the important values are:

\begin{lstlisting}[language=Python]
self.max_angle_vel = 80
self.alpha_fac = -1.0
self.beta_fac = -1.0
\end{lstlisting}

So a policy output near $1.0$ becomes a strong positive or negative velocity command, and a policy output near $0.0$ becomes little or no movement.

\begin{center}
\begin{tabular}{lll}
\toprule
Policy output & Approximate command & Meaning \\
\midrule
$1.0$ & about $\pm 80$ & strong command \\
$0.5$ & about $\pm 40$ & medium command \\
$0.0$ & $0$ & stop or hold \\
$-1.0$ & about $\mp 80$ & strong opposite command \\
\bottomrule
\end{tabular}
\end{center}

\section{Reward and Episode Score}

The reward tells the agent whether it is moving the ball in the correct direction. In this project, reward is mainly based on forward progress along the path:

\begin{lstlisting}[language=Python]
reward = float(curr_pos_path - self.prev_pos_path) * 0.004 / 16.0
\end{lstlisting}

Mathematically:

\begin{equation}
r_t = (p_t - p_{t-1}) \cdot \frac{0.004}{16}
\end{equation}

where:

\begin{itemize}[leftmargin=*]
  \item $p_t$ is the current progress along the path,
  \item $p_{t-1}$ is the previous progress,
  \item $r_t$ is the reward at time $t$.
\end{itemize}

If the ball moves forward, $p_t - p_{t-1}$ is positive, so the reward is positive. If the ball does not move forward, the reward is small. If the ball fails, the episode ends.

The episode score is the sum of reward over the whole episode:

\begin{equation}
R = \sum_{t=0}^{T} r_t
\end{equation}

The score is not a percent. It is a scaled progress score. In the current map, the rough meaning is:

\begin{center}
\begin{tabular}{ll}
\toprule
Score & Approximate meaning \\
\midrule
$0.0$ to $0.2$ & failed early \\
$1.5$ & around half of the map \\
$2.9$ to $3.0$ & near finish or full map \\
\bottomrule
\end{tabular}
\end{center}

\section{Replay Data}

Replay data is saved experience from the robot. Every step produces one training example:

\begin{equation}
(o_t, a_t, r_t, o_{t+1}, d_t)
\end{equation}

where:

\begin{itemize}[leftmargin=*]
  \item $o_t$ is the current observation,
  \item $a_t$ is the action,
  \item $r_t$ is the reward,
  \item $o_{t+1}$ is the next observation,
  \item $d_t$ says whether the episode ended.
\end{itemize}

Dreamer stores replay data in:

\begin{lstlisting}
~/cyberrunner_logs/RUN_NAME/replay/
\end{lstlisting}

This folder contains many \texttt{.npz} files. These files are important because they let Dreamer train from old robot experience. Without replay data, the robot has to collect new experience before it has enough data to learn well.

\section{Why Old Replay Data Matters}

The old run \texttt{robust\_2} contains useful experience. If a new run only loads:

\begin{lstlisting}
robust_2/checkpoint.ckpt
\end{lstlisting}

then it loads the old policy weights, but it may not load the old replay data. That can make startup faster, but it also means the run has less old data available for training.

To continue with both old policy and old data, the new run folder should contain:

\begin{lstlisting}
checkpoint.ckpt
replay/
\end{lstlisting}

This is why the recommended workflow hard-links the old replay folder into a new run folder.

\section{How DreamerV3 Works}

DreamerV3 is a model-based reinforcement learning algorithm. This means it does more than directly memorize which action worked. It learns a model of the environment, then uses that model to improve the policy.

The simplified DreamerV3 loop is:

\begin{enumerate}[leftmargin=*]
  \item Collect real experience from the robot.
  \item Store the experience in replay.
  \item Train a world model from replay.
  \item Use the world model to imagine future trajectories.
  \item Train the policy and value function from imagined trajectories.
  \item Run the improved policy on the robot to collect more data.
\end{enumerate}

The world model learns a compact hidden state:

\begin{equation}
(z_t, a_t) \rightarrow (z_{t+1}, \hat{r}_t)
\end{equation}

where $z_t$ is a latent state. A latent state is a compressed representation of what matters. It is smaller than the raw camera image but still contains useful information for predicting the future.

\subsection{World Model}

The world model tries to predict what will happen next. For CyberRunner, it learns patterns such as:

\begin{itemize}[leftmargin=*]
  \item how the ball moves when the board tilts,
  \item how action affects future ball position,
  \item what kind of state usually leads to reward,
  \item what kind of state usually leads to failure.
\end{itemize}

\subsection{Policy}

The policy chooses the next action. It is trained to choose actions that lead to higher imagined future reward.

\subsection{Value Function}

The value function estimates how good a state is. It answers the question:

\begin{center}
\emph{If the robot is in this state, how much future reward should it expect?}
\end{center}

The policy and value function work together. The policy proposes actions, and the value function helps judge whether those actions lead to good future outcomes.

\section{Why This Is Reinforcement Learning}

This is reinforcement learning because the robot is not given a list of correct servo commands. It must learn from reward.

The objective is:

\begin{equation}
\max_\pi \mathbb{E}\left[\sum_{t=0}^{T} r_t\right]
\end{equation}

In plain language:

\begin{center}
\fbox{\parbox{0.88\linewidth}{
\centering
find a policy that gets the largest total reward by moving the ball farther through the maze
}}
\end{center}

\section{Checkpoints and the Best Policy Problem}

Training saves checkpoints. A normal training run saves:

\begin{lstlisting}
~/cyberrunner_logs/RUN_NAME/checkpoint.ckpt
\end{lstlisting}

This file is usually the latest checkpoint. However, the latest checkpoint is not always the best checkpoint.

For example, the \texttt{robust\_2} score history showed:

\begin{center}
\begin{tabular}{ll}
\toprule
Example & Score \\
\midrule
Best old episode & about $2.93$ \\
Later weak episode & about $0.02$ \\
\bottomrule
\end{tabular}
\end{center}

This means the robot had a very strong episode earlier, but later behavior could still be worse. If only the latest checkpoint is saved, a strong policy can be lost.

It is also important to understand that a score record is not the same as a checkpoint. A line in \texttt{scores.jsonl} tells us what score happened, but it does not contain the policy weights. To reuse a policy, we need the actual checkpoint file from that time.

\section{Training Scripts}

There are three useful training modes.

\subsection{Normal Training}

Normal training uses:

\begin{lstlisting}
--run.script train
\end{lstlisting}

It saves:

\begin{lstlisting}
checkpoint.ckpt
\end{lstlisting}

This is simple, but it only protects the latest checkpoint.

\subsection{Best Checkpoint Training}

The best-checkpoint script is:

\begin{lstlisting}
dreamerv3/dreamerv3/embodied/run/train_best.py
\end{lstlisting}

It is used with:

\begin{lstlisting}
--run.script train_best
\end{lstlisting}

It saves:

\begin{lstlisting}
checkpoint.ckpt
checkpoint_best.ckpt
checkpoint_best_score.txt
\end{lstlisting}

After each episode, it checks whether the episode score is better than the previous best score. If yes, it saves a best checkpoint.

\subsection{Top-10 Checkpoint Training}

The top-10 checkpoint behavior is implemented in:

\begin{lstlisting}
dreamerv3/dreamerv3/embodied/run/train_top5.py
\end{lstlisting}

It is used with:

\begin{lstlisting}
--run.script train_top5
\end{lstlisting}

It saves the top ten checkpoints. Score is the main ranking value. If two
episodes have the same score, the episode with the faster completion time ranks
higher:

\begin{lstlisting}
checkpoint_top_score2.918750_time25.400s_step1458041_episode31.ckpt
checkpoint_top10.txt
checkpoint_top10.json
\end{lstlisting}

This is safer than saving only one checkpoint. One high score can happen by
luck, so keeping ten candidates gives more policies to evaluate later.

\section{Recommended Training Workflow}

The recommended workflow is to continue from the old \texttt{robust\_2} run with both the old checkpoint and old replay data, but use \texttt{train\_top5} so future good policies are saved.

\subsection{Server Command}

Run on the server:

\begin{lstlisting}[language=bash]
cd /home/tbt589/cyberrunner_logs

RUN=robust_2_top5_replay_$(date +%Y%m%d_%H%M%S)
mkdir -p $RUN

cp robust_2/checkpoint.ckpt $RUN/
cp robust_2/config.yaml $RUN/ 2>/dev/null || true
cp robust_2/scores.jsonl $RUN/ 2>/dev/null || true
cp robust_2/metrics.jsonl $RUN/ 2>/dev/null || true
cp -al robust_2/replay $RUN/replay

cd /home/tbt589/cyberruner-main
source install/setup.bash

python3 -m dreamerv3.train \
  --configs cyberrunner large \
  --task gym_cyberrunner_dreamer:cyberrunner-ros-v0 \
  --logdir ~/cyberrunner_logs/$RUN \
  --replay_size 1e6 \
  --run.script train_top5 \
  --run.train_ratio 128 \
  --run.save_every 20 \
  --run.log_every 10 \
  --jax.policy_devices 0 \
  --jax.train_devices 0
\end{lstlisting}

This command does not need \texttt{--run.from\_checkpoint}, because the new run folder already contains \texttt{checkpoint.ckpt} and \texttt{replay/}.

\subsection{Local PC Bridge Command}

Run on the local PC:

\begin{lstlisting}[language=bash]
cd /home/trungbao/CYBER/cyberruner-main
source install/setup.bash

unset ROS_DISCOVERY_SERVER
unset ROS_SUPER_CLIENT
export ROS_DOMAIN_ID=0
export RMW_IMPLEMENTATION=rmw_fastrtps_cpp

python3 tcp_ros_bridge.py aere-a83514.ae.utexas.edu 5555
\end{lstlisting}

\subsection{Checking Saved Top-10 Policies}

To check the current top-ten list:

\begin{lstlisting}[language=bash]
cat ~/cyberrunner_logs/RUN_NAME/checkpoint_top10.txt
\end{lstlisting}

To watch it live:

\begin{lstlisting}[language=bash]
watch -n 5 'cat ~/cyberrunner_logs/RUN_NAME/checkpoint_top10.txt'
\end{lstlisting}

\section{Evaluation}

Evaluation means running a checkpoint without updating the policy. Training changes the policy. Evaluation only tests it.

\begin{center}
\begin{tabular}{ll}
\toprule
Mode & Behavior \\
\midrule
Training & learns, explores, updates policy \\
Evaluation & runs saved policy, no learning \\
\bottomrule
\end{tabular}
\end{center}

Evaluation is important because a single high score may be lucky. A good checkpoint should be tested several times. The best checkpoint is not always the one with the single highest score. The best checkpoint is usually the one with the most reliable success rate.

\section{Why Training Is Difficult}

CyberRunner is difficult because the robot is a real physical system. The agent must learn under noise, delay, and unstable dynamics.

Main difficulties include:

\begin{itemize}[leftmargin=*]
  \item Camera noise can make the ball position inaccurate.
  \item The ball can disappear from the detector.
  \item Servo motion has delay and acceleration limits.
  \item The ball can move faster than the control loop expects.
  \item Small wrong actions can cause failure.
  \item The robot must learn both forward motion and recovery.
  \item Training data is expensive because it comes from the real robot.
  \item The latest checkpoint can be worse than an earlier checkpoint.
  \item A single successful episode can happen by luck.
\end{itemize}

This is why replay data, evaluation, and top checkpoint saving are important.

\section{Practical Lessons}

The most important lessons from this training workflow are:

\begin{itemize}[leftmargin=*]
  \item Saving only \texttt{checkpoint.ckpt} is risky because it stores the latest policy, not necessarily the best policy.
  \item \texttt{scores.jsonl} is useful for analysis, but it does not contain policy weights.
  \item To reuse a good policy, the actual checkpoint file must exist.
  \item Old replay data helps because Dreamer can keep learning from previous robot experience.
  \item \texttt{train\_top5} protects future good policies by saving several high-score checkpoints.
  \item The best policy should be chosen by repeated evaluation, not only by one score.
\end{itemize}

\section{Conclusion}

CyberRunner is trained with reinforcement learning. The DreamerV3 agent observes the robot state, chooses servo velocity actions, receives reward for path progress, and improves its policy using replay data and a learned world model. The policy is saved in checkpoints. Because the latest checkpoint is not always the best, the safest workflow is to continue from a strong old run with replay data and use \texttt{train\_top5} to save future high-score policies. This gives the project a better chance of keeping reliable policies instead of losing them during later training.

\end{document}
